\section{هفته‌ی 15: مدل‌سازی در علوم شناختی 2: یادگیری تقویتی، شناخت بیزی
و هوش مصنوعی}

\begin{quote}
\textbf{فصل 6: نمونه‌هایی از مدل‌سازی در حوزه‌های مختلف علوم شناختی} · درس
1 ساعت · \lr{lab} 2 ساعت \textbf{داده‌ی اصلی:} انتخاب‌های شبیه‌سازی‌شده در
یک تکلیف \lr{bandit} دوبازویی با یک وارونگی \lr{(reversal)}؛ یک آزمایش
شبیه‌سازی‌شده‌ی ترکیب نشانه‌های بینایی\lr{--}لامسه‌ای \lr{(cue combination).}
\end{quote}

\begin{center}\rule{0.5\linewidth}{0.5pt}\end{center}

\subsection{واژه‌های کلیدی}

واژه‌های مهم این درس با معادل فارسی‌شان (حتی اگر در درس‌های قبل آمده
باشند).

\begin{longtable}[]{@{}
  >{\raggedright\arraybackslash}p{(\linewidth - 4\tabcolsep) * \real{0.3333}}
  >{\raggedright\arraybackslash}p{(\linewidth - 4\tabcolsep) * \real{0.3333}}
  >{\raggedright\arraybackslash}p{(\linewidth - 4\tabcolsep) * \real{0.3333}}@{}}
\caption{\centering واژه‌های کلیدی}\tabularnewline
\toprule\noalign{}
\begin{minipage}[b]{\linewidth}\raggedright
\lr{English}
\end{minipage} & \begin{minipage}[b]{\linewidth}\raggedright
فارسی
\end{minipage} & \begin{minipage}[b]{\linewidth}\raggedright
توضیح کوتاه
\end{minipage} \\
\midrule\noalign{}
\endfirsthead
\toprule\noalign{}
\begin{minipage}[b]{\linewidth}\raggedright
\lr{English}
\end{minipage} & \begin{minipage}[b]{\linewidth}\raggedright
فارسی
\end{minipage} & \begin{minipage}[b]{\linewidth}\raggedright
توضیح کوتاه
\end{minipage} \\
\midrule\noalign{}
\endhead
\bottomrule\noalign{}
\endlastfoot
\lr{simulation} & شبیه‌سازی & تولید داده‌ی مصنوعی با اجرای مدل \\
\lr{bandit task} & تکلیف \lr{bandit} & انتخاب مکرر بین گزینه‌هایی با
پاداش احتمالی \\
\lr{cue combination} & ترکیب نشانه‌ها \lr{(cues)} & ترکیب اطلاعات چند حس
برای یک برآورد \\
\lr{Q-learning} & \lr{Q-learning} & قاعده‌ی یادگیری ارزش گزینه‌ها \\
\lr{softmax} / \lr{inverse temperature} (\ensuremath{\beta}) & \lr{softmax} / دمای معکوس
& تبدیل ارزش‌ها به احتمال انتخاب \\
\lr{learning rate} (\ensuremath{\alpha}) & نرخ یادگیری & اندازه‌ی به‌روزرسانی بعد از هر
خطا \\
\lr{temperature} & دما \lr{(temperature)} & پارامتر نویز در انتخاب
\lr{token}؛ \ensuremath{\beta} = \lr{1/T} \\
\lr{trial} & کوشش \lr{(trial)} & یک بار انجام تکلیف توسط شرکت‌کننده \\
\lr{likelihood} & درست‌نمایی & احتمال داده به‌عنوان تابعی از پارامترها \\
\lr{fit} / \lr{fitting} & برازش & تنظیم پارامترها تا مدل با داده جور
شود \\
\lr{parameter recovery} & بازیابی پارامتر & آیا برازش، پارامترهای معلوم
شبیه‌سازی را برمی‌گرداند \\
\lr{parameter} & پارامتر & کمیت قابل تنظیم مدل که بین افراد یا شرایط فرق
می‌کند \\
\lr{ideal observer} & مشاهده‌گر ایده‌آل & مدل استنباط بهینه با اطلاعات
موجود \\
\lr{cognitive model} & مدل شناختی & توصیف دقیق و ریاضی یا محاسباتی از یک
فرایند ذهنی \\
\lr{exploration} / \lr{exploitation} & کاوش / بهره‌برداری & امتحان
گزینه‌های دیگر / انتخاب بهترین گزینه‌ی فعلی \\
\lr{prediction error} & خطای پیش‌بینی & تفاوت پاداش گرفته‌شده و پاداش مورد
انتظار \\
\lr{reinforcement learning (RL)} & یادگیری تقویتی & یادگیری از پاداش و
پیامد \\
\lr{noise} & نویز & تغییرات تصادفی در شواهد یا در داده \\
\lr{estimate} / \lr{estimator} & برآورد / برآوردگر & مقدار به‌دست‌آمده
برای پارامتر / روش به دست آوردن آن \\
\lr{reliability} & پایایی (قابلیت اطمینان) & دقت یک نشانه؛ عکس واریانس
آن \\
\lr{large language model (LLM)} & مدل بزرگ زبانی & شبکه‌ی عصبی‌ای که کلمه‌ی
بعدی متن را پیش‌بینی می‌کند \\
\lr{parameter trade-off} & بده‌بستان پارامترها & تغییر یک پارامتر که با
پارامتر دیگر جبران می‌شود \\
\lr{RLHF} / \lr{reward model} & \lr{RLHF} / مدل پاداش & آموزش
\lr{chatbot} از روی ترجیحات انسانی \\
\lr{Bradley\lr{--}Terry model} & مدل \lr{Bradley\lr{--}Terry} & مدل انتخاب برای
مقایسه‌های جفتی \\
\end{longtable}

\begin{center}\rule{0.5\linewidth}{0.5pt}\end{center}

\subsection{اهداف یادگیری}

در پایان این هفته، دانشجویان می‌توانند:

\begin{enumerate}
\def\labelenumi{\arabic{enumi}.}
\tightlist
\item
  یادگیری از پاداش را با قاعده‌ی \lr{Rescorla\lr{--}Wagner} / \lr{Q-learning}
  و قاعده‌ی انتخاب \lr{softmax} \textbf{صورت‌بندی کنند} و نرخ یادگیری و
  \lr{inverse temperature} را \textbf{تفسیر کنند}.
\item
  درست‌نمایی \lr{trial} به \lr{trial} یک مدل یادگیری را \textbf{بنویسند}،
  آن را با \lr{MLE} \textbf{برازش کنند}، و بازیابی پارامتر و بازیابی مدل
  را در برابر راهبردهای ساده‌تر \textbf{ارزیابی کنند}.
\item
  پیش‌بینی‌های یک مشاهده‌گر ایده‌آل بیزی \lr{(Bayesian ideal observer)} را
  برای ترکیب نشانه‌ها \textbf{استخراج کنند} و \textbf{بیازمایند} که آیا
  رفتار شبیه‌سازی‌شده بهینه است.
\item
  مدل‌های بیزی شناخت (ادعایی درباره‌ی ذهن) را از تحلیل داده‌ی بیزی (روشی
  برای دانشمندان) \textbf{تمیز دهند}.
\item
  درباره‌ی اینکه مدل‌های شناختی و سیستم‌های هوش مصنوعی چطور به هم کمک
  می‌کنند، و اینکه ابزارهای ارزیابی این درس چطور روی سیستم‌های هوش
  مصنوعی‌ای که به‌عنوان مدل شناخت پیشنهاد می‌شوند به کار می‌روند،
  \textbf{بحث کنند}.
\end{enumerate}

\begin{center}\rule{0.5\linewidth}{0.5pt}\end{center}

\subsection{مقدمه}

\textbf{یادگیری از پیامدها.} یک شرکت‌کننده بارها بین دو ماشین بازی انتخاب
می‌کند که یکی بیشتر از دیگری پاداش می‌دهد. وسط کار، ماشین‌ها مخفیانه جایشان
عوض می‌شود. شرکت‌کننده چقدر سریع متوجه می‌شود و خودش را تطبیق می‌دهد؟ چقدر
کاوش \lr{(explore)} می‌کند و چقدر بهره‌برداری \lr{(exploit)}؟ یکی از
موفق‌ترین مدل‌های کل روان‌شناسی و علوم اعصاب با یک معادله جواب می‌دهد:
انتظارتان را به اندازه‌ی کسری از \textbf{خطای پیش‌بینی} \lr{(prediction
error)}، یعنی تفاوت چیزی که گرفتید و چیزی که انتظار داشتید، به‌روز کنید.
همین قاعده‌ی خطامحور در شرطی‌سازی کلاسیک \lr{(Rescorla\lr{--}Wagner)}، در
یادگیری پیوندگرا (قاعده‌ی دلتا)، در شلیک نورون‌های دوپامینی، و در
الگوریتم‌های یادگیری تقویتی پشت سیستم‌های هوش مصنوعی امروزی دیده می‌شود.
برازش آن روی داده‌ی انتخاب یکی از رایج‌ترین کاربردهای مدل‌سازی شناختی امروز
است، مثلاً در روان‌پزشکی محاسباتی.

\textbf{ذهن به‌عنوان آماردان.} وقتی یک شیء را هم‌زمان می‌بینید و لمس
می‌کنید، بینایی و لامسه هر دو برآوردی پرنویز از اندازه‌اش می‌دهند. آیا
آدم‌ها این کار را می‌کنند؟ این کاربردی متفاوت از کلمه‌ی «بیزی» نسبت به
هفته‌های 9\lr{--}13 است: اینجا قاعده‌ی بیز یک \textbf{فرضیه درباره‌ی ذهن} است،
در سطح محاسباتی \lr{Marr} (هفته‌ی 1).

\textbf{مدل‌های شناختی \ensuremath{\leftrightarrow} هوش مصنوعی.} خیلی از ایده‌ها در هر دو جهت بین
مدل‌سازی شناختی و هوش مصنوعی رفت‌وآمد می‌کنند: یادگیری خطامحور، شبکه‌های
عصبی، یادگیری تقویتی، استنباط بیزی. امروز خود سیستم‌های هوش مصنوعی، از
جمله مدل‌های بزرگ زبانی \lr{(LLM}ها)، به‌عنوان مدل شناخت انسان پیشنهاد
می‌شوند. ابزارهای این درس (درست‌نمایی، مقایسه‌ی مدل، آزمون‌های تعمیم،
بازیابی) دقیقاً همان چیزی‌اند که برای ارزیابی چنین ادعاهایی لازم است.

\paragraph{این موضوع کجا به کار می‌رود؟}

\begin{itemize}
\tightlist
\item
  \textbf{روان‌پزشکی محاسباتی:} مدل‌های یادگیری تقویتی روی انتخاب‌های
  بیماران افسرده، معتاد، مبتلا به اسکیزوفرنی یا وسواس برازش می‌شوند تا
  آزموده شود که آیا از پاداش و تنبیه متفاوت یاد می‌گیرند (مثلاً نرخ
  یادگیری پایین‌تر بعد از ضرر).
\item
  \textbf{علوم اعصاب:} خطای پیش‌بینی در قاعده‌ی دلتا با فعالیت نورون‌های
  دوپامینی جور درمی‌آید؛ یکی از شناخته‌شده‌ترین پیوندهای یک مدل شناختی با
  مغز.
\item
  \textbf{آموزش و بازی:} سیستم‌های یادگیری تطبیقی و طراحی بازی از ایده‌های
  \lr{RL} برای مدل کردن اینکه بازیکن‌ها و یادگیرنده‌ها چطور خودشان را با
  بازخورد تطبیق می‌دهند استفاده می‌کنند.
\item
  \textbf{ادراک و رباتیک:} ترکیب بهینه‌ی نشانه‌ها توضیح می‌دهد آدم‌ها چطور
  بینایی، لامسه و شنوایی را ترکیب می‌کنند، و همین قاعده‌ی وزن‌دهی در ادغام
  حسگرها \lr{(sensor fusion)} به کار می‌رود.
\item
  \textbf{هوش مصنوعی:} یادگیری تقویتی امروزی و آموزش \lr{chatbot}ها از
  روی ترجیحات انسانی \lr{(RLHF)} از همین مدل‌های انتخاب استفاده می‌کنند
  (بخش هوش مصنوعی و \lr{LLM}ها را ببینید).
\end{itemize}

\begin{center}\rule{0.5\linewidth}{0.5pt}\end{center}

\Needspace{11\baselineskip}

\subsection{جزوه‌ی درس}

\paragraph{به زبان ساده}

بعد از انتخاب یک گزینه و دیدن نتیجه، ارزش کمی به سمت چیزی که واقعاً اتفاق
افتاد جابه‌جا می‌شود: \emph{ارزش جدید = ارزش قبلی + نرخ یادگیری × (پاداش \ensuremath{-}
ارزش قبلی)}. عبارت داخل پرانتز \textbf{خطای پیش‌بینی} است، یعنی میزان
غافلگیری. \textbf{نرخ یادگیری} \(\alpha\) می‌گوید این جابه‌جایی چقدر بزرگ
است. انتخاب‌ها کاملاً حریصانه نیستند: قاعده‌ی \textbf{\lr{softmax}} گزینه‌ی
باارزش‌تر را بیشتر انتخاب می‌کند و \textbf{\lr{inverse temperature}}
\(\beta\) تعیین می‌کند چقدر بیشتر (\(\beta\) پایین = کاوش تصادفی‌تر،
\(\beta\) بالا = تقریباً همیشه بهترین). برازش این مدل یعنی پیدا کردن
\(\alpha\) و \(\beta\)ای که دنباله‌ی واقعی انتخاب‌های شرکت‌کننده را
محتمل‌ترین کند. \textbf{ذهن بیزی:} وقتی دو حس تخمین‌های پرنویزی از یک چیز
می‌دهند (دیدن و لمس کردن اندازه‌ی یک شیء)، بهترین حدس یک میانگین وزنی است
که به حس قابل اعتمادتر بیشتر اعتماد می‌کند. آدم‌ها اغلب به این وزن‌دهی
بهینه نزدیک‌اند؛ مدل این را بدون هیچ پارامتر آزادی پیش‌بینی می‌کند و همین
آن را یک آزمون قوی می‌کند.

\textbf{زمان‌بندی پیشنهادی (60 دقیقه):} یادگیری از پاداش 20 · برازش
مدل‌های یادگیری 10 · مدل‌های بیزی شناخت 15 · مدل‌های شناختی و هوش مصنوعی 10
· جمع‌بندی درس 5.

\subsubsection{\texorpdfstring{1. یادگیری تقویتی: قاعده‌ی دلتا و
\lr{Q-learning}}{1. یادگیری تقویتی: قاعده‌ی دلتا و Q-learning}}

\textbf{\lr{Rescorla\lr{--}Wagner} / قاعده‌ی دلتا.} بعد از انتخاب گزینه‌ی \(o\)
در \lr{trial} \(t\) و گرفتن پاداش \(r_t \in \{0, 1\}\):

\[
\delta_t = r_t - Q_t(o) \quad\text{\lr{(prediction error)}}, \qquad
Q_{t+1}(o) = Q_t(o) + \alpha\,\delta_t,
\]

و ارزش گزینه‌های انتخاب‌نشده ثابت می‌ماند. \textbf{نرخ یادگیری}
\(\alpha \in [0, 1]\) تعیین می‌کند پیامدهای تازه چقدر قوی انتظارات قدیمی
را بازنویسی کنند. \(\alpha\) بالا سریع تطبیق می‌دهد ولی پرنویز است؛
\(\alpha\) پایین پایدار است ولی وارونگی را دیر می‌فهمد.

\textbf{قاعده‌ی انتخاب \lr{softmax.}} ارزش‌ها با \textbf{\lr{inverse
temperature}} \(\beta \ge 0\) به احتمال انتخاب تبدیل می‌شوند:

\[
P(\text{\lr{choose }} o \text{\lr{ on trial }} t) = \frac{\exp(\beta\, Q_t(o))}{\sum_{o'} \exp(\beta\, Q_t(o'))}.
\]

\(\beta = 0\) یعنی انتخاب تصادفی (کاوش محض)؛ \(\beta\) بزرگ یعنی همیشه
گزینه‌ای که فعلاً بهترین است انتخاب شود (بهره‌برداری). برای دو گزینه این یک
تابع لجستیک از تفاوت ارزش‌هاست:
\(P(o_1) = 1/(1 + e^{-\beta(Q_1 - Q_2)})\).

\textbf{ارتباط‌ها:} - این به‌روزرسانی همان \textbf{قاعده‌ی دلتا}ی شبکه‌های
پیوندگراست \lr{(Verguts,} 2022): یک شبکه‌ی تک‌لایه که با تصحیح خطا آموزش
می‌بیند. - خطاهای پیش‌بینی \(\delta_t\) شبیه پاسخ‌های فازیک نورون‌های
دوپامینی مغز میانی‌اند، و همین مدل‌های \lr{RL} را در علوم اعصاب محاسباتی
محوری کرد. - \lr{Q-learning} با خطاهای \lr{temporal-difference} پایه‌ی
یادگیری تقویتی مدرن در هوش مصنوعی است.

\paragraph{خواندن ریاضی}

\begin{itemize}
\tightlist
\item
  \(Q_t(o)\): ارزش گزینه‌ی \(o\) در \lr{trial} شماره‌ی \(t\) (زیرنویس
  \(t\) شماره‌ی \lr{trial} است).
\item
  \(\delta_t = r_t - Q_t(o)\): پاداش گرفته‌شده منهای پاداش مورد انتظار.
  مثبت = غافلگیری خوشایند، منفی = ناامیدی.
\item
  \(Q_{t+1}(o) = Q_t(o) + \alpha\,\delta_t\). \textbf{مثال:}
  \(Q = 0.5\)، \(\alpha = 0.3\) و پاداش 1: \(\delta = 0.5\) و \(Q\) جدید
  \(= 0.5 + 0.15 = 0.65\). بدون پاداش: \(\delta = -0.5\) و \(Q\) جدید
  \(= 0.35\).
\item
  \textbf{\lr{softmax} برای دو گزینه:}
  \(P(\text{\lr{option 1}}) = 1/(1 + e^{-\beta(Q_1 - Q_2)})\). با
  \(Q_1 - Q_2 = 0.25\): \(\beta = 0\) می‌دهد 0.50 (تصادفی)، \(\beta = 3\)
  می‌دهد 0.68، و \(\beta = 20\) می‌دهد 0.99 (تقریباً همیشه گزینه‌ی 1).
\end{itemize}

\subsubsection{2. برازش یک مدل یادگیری}

\textbf{درست‌نمایی ترتیبی است.} ارزش‌ها به انتخاب‌ها و پاداش‌های گذشته‌ی خود
شرکت‌کننده بستگی دارند، پس مدل را روی دنباله‌ی واقعی \lr{trial}های
شرکت‌کننده اجرا می‌کنیم و احتمال‌هایی را که به انتخاب‌های واقعاً انجام‌شده
نسبت داده در هم ضرب می‌کنیم:

\[
\ell(\alpha, \beta) = \sum_{t=1}^{T} \ln P\big(\text{\lr{choice}}_t \mid Q_t(\cdot;\ \alpha, \text{\lr{history}}_{<t}),\ \beta\big).
\]

\textbf{نکته‌های عملی:} - \lr{reparameterize} کنید:
\(\alpha = 1/(1 + e^{-\eta_\alpha})\) و \(\beta = e^{\eta_\beta}\)
(هفته‌ی 3)، و از \lr{multi-start} استفاده کنید. - \textbf{\lr{trade-off}
\(\alpha\)\lr{--}\(\beta\):} انتخاب پرنویز (\(\beta\) پایین) و یادگیری کند یا
نامنظم (\(\alpha\) حدی) می‌توانند رفتارهای شبیه به هم بسازند، پس برآوردها
در برازش‌های مختلف اغلب همبسته‌اند، و علامت همبستگی به طرح و دامنه‌ی
پارامترها بستگی دارد. مشکل دیگر این است که برآورد مقادیر خیلی بزرگ
\(\beta\) اصلاً سخت است (بخش \lr{C lab).} مطالعه‌ی بازیابی در همان طرح
واقعی ضروری است (هفته‌ی 7). - \textbf{طرح مهم است:} برای برآورد
\(\alpha\) به وارونگی و \lr{trial} کافی نیاز است؛ احتمال‌های پاداش ثابت
اطلاعات کمی درباره‌اش می‌دهند. - \textbf{راهبردهای جایگزین} باید مقایسه
شوند: مثلاً \emph{\lr{win-stay/lose-shift}} \lr{(WSLS:} انتخاب پاداش‌گرفته
را تکرار کن، بعد از بی‌پاداشی عوض کن، با نرخ لغزش \(\varepsilon\))، یا یک
مدل تصادفی سوگیر. برازش خوب \lr{RL} معنای زیادی ندارد اگر \lr{WSLS} هم
به همان خوبی برازش شود.

\paragraph{خواندن ریاضی}

\(\ell(\alpha, \beta) = \sum_t \ln P(\text{\lr{choice}}_t \mid Q_t, \beta)\):
مدل یادگیری را روی دنباله‌ی انتخاب‌ها و پاداش‌های \emph{خود} شرکت‌کننده،
\lr{trial} به \lr{trial}، اجرا کن. در هر \lr{trial} احتمالی را که مدل به
انتخابِ واقعی شرکت‌کننده داده حساب کن، لگاریتمش را بگیر و جمع بزن. مثال
حل‌شده (بخش 1) این کار را برای چهار \lr{trial} انجام می‌دهد:
\(\ln 0.5 + \ln 0.679 + \ln 0.593 + \ln 0.755 = -1.884\). بهترین
\((\alpha, \beta)\) این جمع را بیشینه می‌کنند. این همان درست‌نمایی هفته‌ی 4
است، با این تفاوت که احتمال هر \lr{trial} به تاریخچه‌ی پیش از آن بستگی
دارد.

\subsubsection{3. مدل‌های بیزی شناخت}

\textbf{ترکیب نشانه‌ها.} یک نشانه‌ی بینایی
\(x_V \sim \mathcal{N}(S, \sigma_V)\) و یک نشانه‌ی لامسه‌ای
\(x_H \sim \mathcal{N}(S, \sigma_H)\) هر دو اندازه‌ی واقعی \(S\) را
می‌سنجند. با پیشین تخت، پسین \(S\) نرمال است با میانگین

\[
\hat S = w_V x_V + w_H x_H, \qquad w_V = \frac{1/\sigma_V^2}{1/\sigma_V^2 + 1/\sigma_H^2}, \quad w_H = 1 - w_V,
\]

و انحراف معیار

\[
\sigma_{VH} = \sqrt{\frac{\sigma_V^2\,\sigma_H^2}{\sigma_V^2 + \sigma_H^2}} \;\le\; \min(\sigma_V, \sigma_H).
\]

\textbf{پیش‌بینی‌های آزمون‌پذیر} (مثلاً آزموده‌شده توسط \lr{Ernst} \lr{\&}
\lr{Banks,} 2002): (1) وزن‌ها باید پایایی نشانه‌ها را دنبال کنند، که
آزمایشگر می‌تواند با تخریب یک نشانه دست‌کاری‌اش کند؛ (2) دقت دوحسی باید
\emph{بهتر} از بهترین نشانه‌ی تکی باشد، دقیقاً به اندازه‌ی بالا. انحراف
معیارهای تک‌نشانه جداگانه اندازه‌گیری می‌شوند، پس پیش‌بینی‌های مشاهده‌گر
ایده‌آل \textbf{هیچ پارامتر آزادی ندارند}. همین آن‌ها را آزمون‌هایی قوی و
ابطال‌پذیر می‌کند.

\textbf{یک مدل بیزی شناخت چه ادعایی دارد.} اینکه رفتار، با توجه به نویز،
پیشین و هدف‌های عامل، تقریبی از استنباط بهینه است: یک تبیین در
\textbf{سطح محاسباتی} \lr{(Marr).}

\textbf{نقدهایی که باید در ذهن داشت:} با انتخاب آزادانه‌ی پیشین‌ها،
درست‌نمایی‌ها و تابع‌های زیان، تقریباً هر رفتاری را می‌شود بعد از واقعه
«بهینه» نامید (مسئله‌ی انعطاف‌پذیری هفته‌ی 6). آزمون‌های قوی این اجزا را
مستقل تثبیت می‌کنند، مثل ترکیب نشانه‌ها.

\textbf{دو معنا، یک بار دیگر:}

\begin{longtable}[]{@{}
  >{\raggedright\arraybackslash}p{(\linewidth - 4\tabcolsep) * \real{0.3333}}
  >{\raggedright\arraybackslash}p{(\linewidth - 4\tabcolsep) * \real{0.3333}}
  >{\raggedright\arraybackslash}p{(\linewidth - 4\tabcolsep) * \real{0.3333}}@{}}
\caption{\centering 3. مدل‌های بیزی شناخت}\tabularnewline
\toprule\noalign{}
\begin{minipage}[b]{\linewidth}\raggedright
\end{minipage} & \begin{minipage}[b]{\linewidth}\raggedright
مدل بیزی شناخت
\end{minipage} & \begin{minipage}[b]{\linewidth}\raggedright
تحلیل داده‌ی بیزی
\end{minipage} \\
\midrule\noalign{}
\endfirsthead
\toprule\noalign{}
\begin{minipage}[b]{\linewidth}\raggedright
\end{minipage} & \begin{minipage}[b]{\linewidth}\raggedright
مدل بیزی شناخت
\end{minipage} & \begin{minipage}[b]{\linewidth}\raggedright
تحلیل داده‌ی بیزی
\end{minipage} \\
\midrule\noalign{}
\endhead
\bottomrule\noalign{}
\endlastfoot
چه کسی استنباط می‌کند؟ & ذهن شرکت‌کننده (به‌طور فرضی) & دانشمند \\
پیشین نشان‌دهنده‌ی & باورهای شرکت‌کننده درباره‌ی جهان & دانش دانشمند درباره‌ی
پارامترها \\
می‌تواند غلط باشد اگر\ensuremath{\ldots} & آدم‌ها بهینه نباشند & مدل یا پیشین‌ها بد
انتخاب شده باشند \\
نمونه‌ها در این درس & ترکیب نشانه‌ها (امروز) & هفته‌های 9\lr{--}13 \\
\end{longtable}

می‌توانید یک مدل بیزی شناخت را با بیشینه‌ی درست‌نمایی تحلیل کنید، و یک مدل
غیربیزی (مثلاً \lr{Q-learning)} را با تحلیل داده‌ی بیزی.

\paragraph{خواندن ریاضی}

\begin{itemize}
\tightlist
\item
  \(w_V = \dfrac{1/\sigma_V^2}{1/\sigma_V^2 + 1/\sigma_H^2}\): وزن
  بینایی برابر دقت آن (1/واریانس) تقسیم بر دقت کل است. حسی که دو برابر
  دقیق‌تر باشد (نصف انحراف معیار، یک‌چهارم واریانس) چهار برابر وزن می‌گیرد.
\item
  \(\sigma_{VH} = \sqrt{\sigma_V^2\sigma_H^2/(\sigma_V^2 + \sigma_H^2)}\):
  انحراف معیار ترکیبی که همیشه از بهترین حس تنها کوچک‌تر است.
  \textbf{مثال:} \(\sigma_V = 1\) و \(\sigma_H = 2\):
  \(w_V = 1/(1 + 1/4) = 0.8\) و \(\sigma_{VH} = \sqrt{4/5} = 0.894\)،
  حدود 11\% بهتر از بینایی به‌تنهایی.
\item
  \textbf{آزمون:} \(\sigma_V\) و \(\sigma_H\) را جداگانه اندازه بگیر (هر
  بار یک حس)، \(w_V\) و \(\sigma_{VH}\) را بدون هیچ پارامتر آزادی
  پیش‌بینی کن، و بعد داده‌ی دوحسی را بررسی کن.
\end{itemize}

\paragraph{برای یک مهندس}

\begin{itemize}
\tightlist
\item
  \textbf{قاعده‌ی دلتا همان به‌روزرسانی \lr{LMS} است.}
  \(Q \leftarrow Q + \alpha(r - Q)\) یک \lr{LMS} تک‌ضریبی است (یا یک
  میانگین متحرک با وزن نمایی و ضریب فراموشی \(1-\alpha\)): \(\alpha\)
  بزرگ تغییرات را سریع دنبال می‌کند ولی پرنویز است، \(\alpha\) کوچک هموار
  ولی کند است؛ همان \lr{trade-off} آشنای ردیابی و خطای تنظیم
  \lr{(misadjustment).}
\item
  \textbf{\lr{softmax}} همان توزیع \lr{Boltzmann/Gibbs} است؛ \(\beta\)
  به معنای فیزیکی یک دمای معکوس است. نمونه‌گیری با دما همان روشی است که
  مدل‌های زبانی متن تولید می‌کنند (کادر هوش مصنوعی هفته‌ی 1).
\item
  \textbf{برازش \(\alpha\) و \(\beta\) از روی انتخاب‌ها} یعنی شناسایی یک
  سیستم تطبیقی فقط از روی خروجی‌های دودویی آن: هرگز \(Q\) را نمی‌بینید،
  فقط تصمیم‌هایی را که تولید می‌کند. به همین دلیل \(\alpha\) و \(\beta\)
  با هم \lr{trade-off} دارند و بررسی بازیابی ضروری است.
\item
  \textbf{ترکیب بهینه‌ی نشانه‌ها} دقیقاً از وزن‌های معکوس واریانس در
  \textbf{ترکیب با نسبت بیشینه \lr{(MRC)}} و ادغام حسگرها در فیلتر کالمن
  استفاده می‌کند: واریانس تخمین ترکیبی
  \((1/\sigma_V^2 + 1/\sigma_H^2)^{-1}\) است. سؤال علمی این است که آیا
  \emph{مغز} این ادغام را انجام می‌دهد.
\end{itemize}

\subsubsection{4. مدل‌های شناختی و هوش مصنوعی}

\begin{longtable}[]{@{}
  >{\raggedright\arraybackslash}p{(\linewidth - 2\tabcolsep) * \real{0.5000}}
  >{\raggedright\arraybackslash}p{(\linewidth - 2\tabcolsep) * \real{0.5000}}@{}}
\caption{\centering 4. مدل‌های شناختی و هوش مصنوعی}\tabularnewline
\toprule\noalign{}
\begin{minipage}[b]{\linewidth}\raggedright
جهت
\end{minipage} & \begin{minipage}[b]{\linewidth}\raggedright
نمونه‌ها
\end{minipage} \\
\midrule\noalign{}
\endfirsthead
\toprule\noalign{}
\begin{minipage}[b]{\linewidth}\raggedright
جهت
\end{minipage} & \begin{minipage}[b]{\linewidth}\raggedright
نمونه‌ها
\end{minipage} \\
\midrule\noalign{}
\endhead
\bottomrule\noalign{}
\endlastfoot
علوم شناختی \ensuremath{\leftarrow} هوش مصنوعی & شبکه‌های عصبی (الهام‌گرفته از نورون‌ها و
قاعده‌های یادگیری)؛ یادگیری تقویتی (از نظریه‌ی یادگیری حیوانی)؛ ایده‌های
\lr{exemplar/prototype} در طبقه‌بندهای یادگیری ماشین (روش‌های نزدیک‌ترین
همسایه در برابر مرکزوار) \\
هوش مصنوعی \ensuremath{\leftarrow} علوم شناختی & شبکه‌های عمیق به‌عنوان مدل قشر بینایی؛ عامل‌های
\lr{RL} به‌عنوان مدل تصمیم‌گیری؛ مدل‌های بزرگ زبانی که به‌عنوان مدل پردازش
زبان و استدلال انسان‌گونه پیشنهاد می‌شوند \\
ریاضیات مشترک & درست‌نمایی‌ها و تابع‌های زیان؛ بهینه‌سازی مبتنی بر گرادیان
(هفته‌ی 3)؛ استنباط بیزی و \lr{MCMC}؛ \lr{cross-validation} و تعمیم
(هفته‌ی 6) \\
\end{longtable}

\textbf{ارزیابی یک سیستم هوش مصنوعی به‌عنوان مدل شناختی} دقیقاً از
جعبه‌ابزار درس استفاده می‌کند: - آیا داده‌ی انسانی را \emph{فراتر} از برازش
پیش‌بینی می‌کند، مثلاً شرایط تازه و تفاوت‌های فردی را (تعمیم، هفته‌ی 6)؟ -
آیا در یک مقایسه‌ی منصفانه از مدل‌های شناختی ساده‌تر و تفسیرپذیر بهتر است
(هفته‌های 7، 10)؟ - آیا می‌توانیم از روی داده بگوییم انسان و مدل از یک
سازوکار استفاده می‌کنند، یا فقط خروجی‌های شبیه تولید می‌کنند (تقلید مدل و
بازیابی، هفته‌ی 7)؟ - آیا پارامترها یا حالت‌های درونی‌اش به‌عنوان سازه‌های
روان‌شناختی تفسیرپذیرند (هفته‌ی 1)؟

مدلی با میلیاردها پارامتر که رفتار را خوب برازش می‌کند به‌طور خودکار یک
\emph{تبیین} نیست: انعطاف‌پذیری، \lr{identifiability} و تفسیرپذیری به
اندازه‌ی برازش مهم‌اند.

\subsubsection{5. جمع‌بندی درس}

در حوزه‌های خیلی متفاوت (فراموشی، تشخیص، دسته‌بندی، تصمیم‌گیری، یادگیری،
ادراک)، گردش کار یکسان بوده است: \textbf{صورت‌بندی \ensuremath{\leftarrow} شبیه‌سازی \ensuremath{\leftarrow} برازش \ensuremath{\leftarrow}
ارزیابی \ensuremath{\leftarrow} بازیابی \ensuremath{\leftarrow} تفسیر.} حوزه مدل را عوض می‌کند؛ منطق مدل‌سازی ثابت
می‌ماند.

\paragraph{برای مدرس}

\begin{itemize}
\tightlist
\item
  \textbf{بازی \lr{bandit} را با کلاس اجرا کنید.} دو «دستگاه» (دو دانشجو
  با سکه‌ی پنهان یا یک اپلیکیشن عدد تصادفی؛ یکی 80\% مواقع پاداش می‌دهد و
  دیگری 20\%)، و در وسط کار مخفیانه جایشان را عوض کنید. کلاس پیش از دیدن
  هر معادله‌ای کاوش، بهره‌برداری و وارونگی را تجربه می‌کند.
\item
  \textbf{چهار \lr{trial Q-learning} مثال حل‌شده را روی تخته انجام دهید}
  با یک جدول: ارزش قبل، انتخاب، پاداش، خطای پیش‌بینی، ارزش بعد.
\item
  \textbf{برای ترکیب نشانه‌ها آزادانه از شهود مهندسی‌تان استفاده کنید:}
  «ترکیب دو حسگر پرنویز» راه طبیعی ارائه‌ی آن است و برای دانشجوهای
  روان‌شناسی روشنگر است.
\item
  \textbf{بحث هوش مصنوعی را روی زمین نگه دارید.} بخش 4 می‌تواند بحث‌های
  طولانی راه بیندازد؛ از چهار سؤال ارزیابی (تعمیم، مقایسه‌ی منصفانه،
  تقلید، تفسیرپذیری) برای متمرکز نگه داشتن آن استفاده کنید.
\item
  \textbf{سؤال‌های رایج.}

  \begin{itemize}
  \tightlist
  \item
    \emph{«چرا برآورد \(\beta\) سخت است؟»} وقتی انتخاب‌ها تقریباً قطعی
    شدند، بیشتر کردن \(\beta\) درست‌نمایی را تقریباً تغییر نمی‌دهد.
  \item
    \emph{«آیا مغز واقعاً \lr{Q-learning} انجام می‌دهد؟»} توصیف خوبی از
    بسیاری از داده‌های انتخاب است و با سیگنال‌های دوپامین جور درمی‌آید، ولی
    مدل‌های پیچیده‌تر (یادگیری مبتنی بر مدل یا \lr{model-based)} هم
    پشتیبانی دارند.
  \item
    \emph{«آیا آدم‌ها همیشه نشانه‌ها را بهینه ترکیب می‌کنند؟»} اغلب نزدیک،
    ولی نه همیشه؛ شرکت‌کننده‌ی \lr{lab} به لامسه بیش از حد وزن می‌دهد.
  \end{itemize}
\end{itemize}

\begin{center}\rule{0.5\linewidth}{0.5pt}\end{center}

\Needspace{11\baselineskip}

\subsection{مثال حل‌شده}

\subsubsection{\texorpdfstring{بخش 1: چهار \lr{trial Q-learning} با
دست}{بخش 1: چهار trial Q-learning با دست}}

دو گزینه؛ \(\alpha = 0.5\)، \(\beta = 3\)؛ ارزش‌های اولیه
\(Q(1) = Q(2) = 0.5\).

\begin{landscape}\small
\begin{longtable}[]{@{}
  >{\raggedright\arraybackslash}p{(\linewidth - 14\tabcolsep) * \real{0.1250}}
  >{\raggedright\arraybackslash}p{(\linewidth - 14\tabcolsep) * \real{0.1250}}
  >{\raggedright\arraybackslash}p{(\linewidth - 14\tabcolsep) * \real{0.1250}}
  >{\raggedright\arraybackslash}p{(\linewidth - 14\tabcolsep) * \real{0.1250}}
  >{\raggedright\arraybackslash}p{(\linewidth - 14\tabcolsep) * \real{0.1250}}
  >{\raggedright\arraybackslash}p{(\linewidth - 14\tabcolsep) * \real{0.1250}}
  >{\raggedright\arraybackslash}p{(\linewidth - 14\tabcolsep) * \real{0.1250}}
  >{\raggedright\arraybackslash}p{(\linewidth - 14\tabcolsep) * \real{0.1250}}@{}}
\caption{\centering بخش 1: چهار \lr{trial Q-learning }با دست}\tabularnewline
\toprule\noalign{}
\begin{minipage}[b]{\linewidth}\raggedright
\lr{trial}
\end{minipage} & \begin{minipage}[b]{\linewidth}\raggedright
\(Q(1)\)
\end{minipage} & \begin{minipage}[b]{\linewidth}\raggedright
\(Q(2)\)
\end{minipage} & \begin{minipage}[b]{\linewidth}\raggedright
\(P(\text{\lr{choose }} 1)\)
\end{minipage} & \begin{minipage}[b]{\linewidth}\raggedright
انتخاب
\end{minipage} & \begin{minipage}[b]{\linewidth}\raggedright
پاداش
\end{minipage} & \begin{minipage}[b]{\linewidth}\raggedright
\(\delta\)
\end{minipage} & \begin{minipage}[b]{\linewidth}\raggedright
به‌روزرسانی
\end{minipage} \\
\midrule\noalign{}
\endfirsthead
\toprule\noalign{}
\begin{minipage}[b]{\linewidth}\raggedright
\lr{trial}
\end{minipage} & \begin{minipage}[b]{\linewidth}\raggedright
\(Q(1)\)
\end{minipage} & \begin{minipage}[b]{\linewidth}\raggedright
\(Q(2)\)
\end{minipage} & \begin{minipage}[b]{\linewidth}\raggedright
\(P(\text{\lr{choose }} 1)\)
\end{minipage} & \begin{minipage}[b]{\linewidth}\raggedright
انتخاب
\end{minipage} & \begin{minipage}[b]{\linewidth}\raggedright
پاداش
\end{minipage} & \begin{minipage}[b]{\linewidth}\raggedright
\(\delta\)
\end{minipage} & \begin{minipage}[b]{\linewidth}\raggedright
به‌روزرسانی
\end{minipage} \\
\midrule\noalign{}
\endhead
\bottomrule\noalign{}
\endlastfoot
1 & 0.50 & 0.50 & 0.500 & 1 & 1 & \(1 - 0.5 = +0.5\) &
\(Q(1) = 0.5 + 0.5 \times 0.5 = 0.75\) \\
2 & 0.75 & 0.50 & \(1/(1 + e^{-3 \times 0.25}) = 0.679\) & 1 & 0 &
\(0 - 0.75 = -0.75\) & \(Q(1) = 0.75 - 0.375 = 0.375\) \\
3 & 0.375 & 0.50 & \(1/(1 + e^{-3 \times (-0.125)}) = 0.407\) & 2 & 1 &
\(1 - 0.5 = +0.5\) & \(Q(2) = 0.75\) \\
4 & 0.375 & 0.75 & \(1/(1 + e^{-3 \times (-0.375)}) = 0.245\) & 2 & 1 &
\(+0.25\) & \(Q(2) = 0.875\) \\
\end{longtable}
\end{landscape}

\textbf{لگاریتم درست‌نمایی این چهار انتخاب:}
\(\ln 0.5 + \ln 0.679 + \ln(1 - 0.407) + \ln(1 - 0.245) = -0.693 - 0.387 - 0.523 - 0.281 = -1.884\).

\textbf{تفسیر.} با \(\alpha = 0.5\) یک بار بی‌پاداشی ارزش گزینه‌ی 1 را نصف
می‌کند و مدل یک تغییر انتخاب پیش‌بینی می‌کند، که شرکت‌کننده انجامش داد.
درست‌نمایی به مقادیر پارامتری پاداش می‌دهد که دنباله‌ی انتخاب واقعی
شرکت‌کننده را پیش‌بینی کنند.

\subsubsection{بخش 2: یک مشاهده‌گر ایده‌آل}

انحراف معیار بینایی \(\sigma_V = 1\) \lr{mm}، انحراف معیار لامسه
\(\sigma_H = 2\) \lr{mm.} وزن‌ها: \(w_V = 1/(1 + 1/4) = 0.8\)،
\(w_H = 0.2\). انحراف معیار ترکیبی پیش‌بینی‌شده:
\(\sqrt{(1 \times 4)/(1 + 4)} = 0.894\) \lr{mm}، حدود 11\% بهتر از
بینایی به تنهایی. اگر \(x_V = 50\) و \(x_H = 54\) \lr{mm} باشد، برآورد
بهینه \(0.8 \times 50 + 0.2 \times 54 = 50.8\) \lr{mm} است.

\textbf{تفسیر.} مشاهده‌گری که نشانه‌ها را فقط \emph{میانگین} بگیرد
(\(w = 0.5\)) انحراف معیار
\(\sqrt{0.25 \times 1 + 0.25 \times 4} = 1.118\) \lr{mm} دارد، یعنی
\emph{بدتر} از بینایی به تنهایی. وزن‌دهی بهینه است که ترکیب نشانه‌ها را
ارزشمند می‌کند.

\begin{center}\rule{0.5\linewidth}{0.5pt}\end{center}

\Needspace{11\baselineskip}

\subsection{\texorpdfstring{\lr{lab} (2 ساعت)}{lab (2 ساعت)}}

\subsubsection{اهداف}

\begin{enumerate}
\def\labelenumi{\arabic{enumi}.}
\tightlist
\item
  شبیه‌سازی یک عامل \lr{Q-learning} در یک تکلیف \lr{bandit} با وارونگی.
\item
  برازش \((\alpha, \beta)\) با \lr{MLE}؛ اجرای بازیابی پارامتر و تشخیص
  \lr{trade-off} \(\alpha\)\lr{--}\(\beta\).
\item
  مقایسه‌ی \lr{Q-learning} با مدل‌های \lr{win-stay/lose-shift} و تصادفی
  سوگیر.
\item
  آزمودن اینکه آیا یک مشاهده‌گر شبیه‌سازی‌شده نشانه‌ها را بهینه ترکیب می‌کند.
\end{enumerate}

\textbf{زمان‌بندی:} بخش \lr{A} 20 دقیقه · بخش \lr{B} 20 دقیقه · بخش
\lr{C} 25 دقیقه · بخش \lr{D} 30 دقیقه · بخش \lr{E} 25 دقیقه.

\subsubsection{\texorpdfstring{بخش \lr{C:} بازیابی پارامتر و
\lr{trade-off} \ensuremath{\alpha}\lr{--}\ensuremath{\beta} (25
دقیقه)}{بخش C: بازیابی پارامتر و trade-off \ensuremath{\alpha}\lr{--}\ensuremath{\beta} (25 دقیقه)}}

\textbf{تفسیر کنید:} با 200 \lr{trial} و یک وارونگی، \(\alpha\) خوب
بازیابی می‌شود \lr{(r} = 0.90) ولی \(\beta\) تقریباً اصلاً نه \lr{(r} =
0.17). این برازش بخش \lr{B} را هم توضیح می‌دهد (\(\hat\beta = 3.8\) برای
مقدار واقعی 5.0). به ناحیه‌ی \(\beta\) بزرگ در نمودار \(\beta\) نگاه
کنید. چرا دماهای معکوس بالا را به سختی می‌شود تعیین کرد؟ (راهنمایی: وقتی
انتخاب‌ها تقریباً قطعی شدند، بالاتر بردن \(\beta\) درست‌نمایی را به‌زحمت
تغییر می‌دهد.) \(\alpha\) و \(\beta\) بازیابی‌شده هم همبسته‌اند (0.31) با
اینکه مقادیر واقعی مستقل بودند: یک معاوضه. این برای مطالعه‌ای که
\(\beta\) را با مثلاً یک پرسش‌نامه‌ی بالینی همبسته می‌کند چه معنایی دارد؟

\subsubsection{\texorpdfstring{بخش \lr{E:} یک مشاهده‌گر ایده‌آل برای ترکیب
نشانه‌ها (25
دقیقه)}{بخش E: یک مشاهده‌گر ایده‌آل برای ترکیب نشانه‌ها (25 دقیقه)}}

یک شرکت‌کننده‌ی شبیه‌سازی‌شده اندازه‌ی اشیایی را که هم‌زمان دیده و لمس می‌شوند
قضاوت می‌کند. سطح نویز تک‌نشانه را از بلوک‌های تک‌حسی می‌دانیم. آیا شرکت‌کننده
نشانه‌ها را بهینه ترکیب می‌کند؟

\textbf{تفسیر کنید:} آیا این شرکت‌کننده بهینه است؟ وزن بینایی تجربی
(0.56) خیلی پایین‌تر از مقدار بهینه‌ی 0.80 است: شرکت‌کننده به نشانه‌ی
لامسه‌ای کم‌پایاتر بیش از حد وزن می‌دهد. در نتیجه انحراف معیار دوحسی‌اش
(1.006 \lr{mm)} بهتر از بینایی به تنهایی (1.000 \lr{mm)} نیست، به جای
0.894 \lr{mm} که یک مشاهده‌گر ایده‌آل به آن می‌رسید. توجه کنید که پیش‌بینی
مشاهده‌گر ایده‌آل \textbf{هیچ پارامتر آزادی} نداشت: همین آن را یک آزمون
قوی می‌کند.

(این شرکت‌کننده را با وزن واقعی 0.65 شبیه‌سازی کردیم. برآورد 0.56 حدود 3
خطای استاندارد فاصله دارد، یک نمونه‌ی غیرمعمول بدشانس ولی یادآوری مفید که
وزن‌های تجربی هم خطای نمونه‌گیری دارند. تمرین 4 بازه‌های \lr{bootstrap}
دورشان می‌گذارد.)

\subsubsection{تمرین‌ها (به ترتیب سختی)}

آیا لگاریتم درست‌نمایی \ensuremath{-}1.884 به دست می‌آید؟

\textbf{تمرین 2 \ensuremath{\star}.} بخش \lr{C} را با 50 \lr{trial} (بدون وارونگی) و با
400 \lr{trial} (سه وارونگی) دوباره اجرا کنید. بازیابی \(\alpha\) چطور
تغییر می‌کند؟ این برای طراحی آزمایش چه معنایی دارد؟

\textbf{تمرین 3 \ensuremath{\star}\ensuremath{\star} (دو نرخ یادگیری).} بسیاری از مطالعه‌ها نرخ‌های یادگیری
متفاوتی برای خطاهای پیش‌بینی مثبت و منفی پیدا می‌کنند. مدلی با
\(\alpha^{+}\) و \(\alpha^{-}\) صورت‌بندی کنید، یک عامل با
\(\alpha^{+} = 0.4\) و \(\alpha^{-} = 0.1\) شبیه‌سازی کنید، و با
\lr{AIC/BIC} و یک مطالعه‌ی بازیابی بیازمایید که آیا پارامتر اضافه در طرح
بخش \lr{A identifiable} است.

\(w_V = 0.8\) · \(w_V = 0.5\) با \lr{bootstrap} (هفته‌ی 5) دور وزن‌ها و
انحراف معیارهای تجربی بازه‌ی اطمینان بگذارید. با 100 \lr{trial} می‌توانید
دو مشاهده‌گر را از هم تشخیص دهید؟

\((\alpha_j, \beta_j)\)

\subsubsection{\texorpdfstring{\lr{lab check-off}}{lab check-off}}

نشان دهید: (1) برازش \lr{MLE} و نمودارهای بازیابی همراه با تفسیرتان از
نمودار \(\beta\)، (2) جدول بازیابی مدل، و (3) جوابتان درباره‌ی بهینه بودن
مشاهده‌گر ترکیب نشانه‌ها.

\begin{center}\rule{0.5\linewidth}{0.5pt}\end{center}

\subsection{\texorpdfstring{ارتباط با هوش مصنوعی و
\lr{LLM}ها}{ارتباط با هوش مصنوعی و LLMها}}

\textbf{قاعده‌ی انتخاب پشت آموزش \lr{chatbot}ها.} \lr{chatbot}های امروزی
با \emph{یادگیری تقویتی از بازخورد انسانی} \lr{(RLHF)} تنظیم دقیق
\lr{(fine-tune)} می‌شوند. مدل پاداش معمولاً مدل
\textbf{\lr{Bradley\lr{--}Terry}} است،
\(P(i \succ j) = 1/(1 + e^{-(r_i - r_j)})\). این همان قاعده‌ی انتخاب
\lr{softmax} این هفته با \(\beta = 1\)، و قاعده‌ی انتخاب \lr{Luce} در
\lr{GCM} است. برازش آن یک مسئله‌ی \lr{MLE} از هفته‌ی 4 است. پس مدل‌های
انتخاب انسانی در قلب شیوه‌ای قرار دارند که سیستم‌های هوش مصنوعی امروزی با
ترجیحات انسانی همسو \lr{(align)} می‌شوند، و مشکلاتشان هم همین‌جاست:
انتخاب‌های پرنویز، \lr{identifiability} (فقط \emph{تفاوت}‌های \(r\)
مهم‌اند، پس یک مقدار ثابت می‌شود)، و فاصله‌ی بین آنچه آدم‌ها انتخاب می‌کنند و
آنچه واقعاً خوب است.

با حدود 300 مقایسه‌ی پرنویز، ترتیب برآوردشده‌ی پاسخ‌های 1 و 2 حتی برعکس
شده. یک مدل پاداش یک مدل شناختی برازش‌شده است، و مثل هر مدل دیگری به
برآورد عدم قطعیت و بررسی بازیابی نیاز دارد.

آیا برعکس شدن پاسخ‌های 1 و 2 در محدوده‌ی عدم قطعیت است؟ - \ensuremath{\star}\ensuremath{\star} یک بازیابی
پارامتر اجرا کنید: چند مقایسه لازم است تا ترتیب هر چهار پاسخ در 90\%
مجموعه داده‌های شبیه‌سازی‌شده درست بازیابی شود؟ - \ensuremath{\star}\ensuremath{\star} ارزیاب‌هایی را شبیه‌سازی
کنید که جواب‌های \emph{طولانی‌تر} را صرف‌نظر از کیفیت ترجیح می‌دهند (یک
پاداش طول به پاداش واقعی اضافه کنید). مدل پاداش برازش‌شده چه یاد می‌گیرد،
و بعد \lr{chatbot} برای انجام چه کاری آموزش می‌بیند؟

\emph{مطالعه‌ی اختیاری:} \lr{Christiano, P. F., et al.~(2017). Deep
reinforcement learning from human preferences.} \emph{\lr{NeurIPS}}.
\lr{Ouyang, L., et al.~(2022). Training language models to follow
instructions with human feedback.} \emph{\lr{NeurIPS}}.

\begin{center}\rule{0.5\linewidth}{0.5pt}\end{center}

\subsection{اشتباهات رایج و برداشت‌های نادرست}

\begin{enumerate}
\def\labelenumi{\arabic{enumi}.}
\tightlist
\item
  \textbf{برازش یک مدل یادگیری روی نسبت‌های انتخاب تجمیع‌شده.} درست‌نمایی
  باید دنباله‌ی \lr{trial}های خود هر شرکت‌کننده را دنبال کند، چون ارزش‌ها
  به تاریخچه‌ی خودش بستگی دارند.
\item
  \textbf{تفسیر \(\alpha\) بدون بررسی بازیابی.} \(\alpha\) و \(\beta\)
  با هم \lr{trade-off} دارند، و \(\alpha\) بدون وارونگی یا \lr{trial}
  کافی بد شناسایی می‌شود.
\item
  \textbf{مقایسه نکردن با راهبردهای ساده.} \lr{win-stay/lose-shift} یا
  پافشاری \lr{(perseveration)} می‌توانند یادگیری را تقلید کنند. نشان دهید
  مدل \lr{RL} در یک مقایسه‌ی منصفانه برنده است.
\item
  \textbf{فراموش کردن اینکه \lr{softmax} از \emph{تفاوت} ارزش‌ها استفاده
  می‌کند.} افزودن یک ثابت به همه‌ی ارزش‌های \lr{Q} چیزی را عوض نمی‌کند، در
  حالی که مقیاس ارزش‌ها و \(\beta\) با هم \lr{trade-off} دارند.
\item
  \textbf{«آدم‌ها بیزی‌اند» از روی یک برازش پس‌از-واقعه.} اگر پیشین و
  درست‌نمایی برای برازش رفتار انتخاب شده باشند، بهینگی فرض شده، نه
  آزموده. آزمون‌های قوی آن‌ها را مستقل تثبیت می‌کنند (ترکیب نشانه‌ها).
\item
  \textbf{قاطی کردن دو نوع بیزی‌گرایی} (مدل ذهن در برابر روش تحلیل).
\item
  \textbf{«یک مدل بزرگ هوش مصنوعی که داده‌ی انسانی را برازش می‌کند، مدل
  شناختی خوبی است.»} برازش تبیین نیست. مثل هر مدل دیگری درباره‌ی تعمیم،
  تقلید، \lr{identifiability} و تفسیرپذیری بپرسید.
\end{enumerate}

\begin{center}\rule{0.5\linewidth}{0.5pt}\end{center}

\subsection{خودآزمایی}

\textbf{1.} به‌روزرسانی قاعده‌ی دلتا را بنویسید و نقش خطای پیش‌بینی را
توضیح دهید. \textbf{پاسخ} \(Q \leftarrow Q + \alpha(r - Q)\). خطای
پیش‌بینی \(r - Q\) تفاوت پیامد و انتظار است: خطاهای مثبت ارزش را بالا
می‌برند، خطاهای منفی پایین می‌آورند، و نرخ یادگیری اندازه‌ی تغییر را تعیین
می‌کند.

\textbf{2.} \lr{inverse temperature} \(\beta\) چه چیزی را کنترل می‌کند؟
\(\beta = 0\) و \(\beta \to \infty\) چه رفتاری تولید می‌کنند؟
\textbf{پاسخ} حساسیت انتخاب‌ها به تفاوت ارزش‌ها (کاوش در برابر
بهره‌برداری). \(\beta = 0\): انتخاب تصادفی بدون توجه به ارزش‌ها.
\(\beta \to \infty\): همیشه گزینه‌ای با بالاترین ارزش فعلی انتخاب می‌شود.

\textbf{3.} چرا درست‌نمایی یک مدل یادگیری باید \lr{trial} به \lr{trial}
روی دنباله‌ی واقعی هر شرکت‌کننده حساب شود؟ \textbf{پاسخ} پیش‌بینی‌های مدل در
\lr{trial} \(t\) به انتخاب‌ها و پاداش‌های قبلی خود آن شرکت‌کننده (از طریق
ارزش‌های \lr{Q)} بستگی دارد، پس احتمال هر انتخاب مشاهده‌شده را فقط با
اجرای مدل روی تاریخچه‌ی مشاهده‌شده می‌شود حساب کرد.

\textbf{4.} \lr{trade-off} \(\alpha\)\lr{--}\(\beta\) را توضیح دهید.
\textbf{پاسخ} ترکیب‌های مختلف نرخ یادگیری و نویز انتخاب می‌توانند
دنباله‌های انتخاب مشابهی تولید کنند. مثلاً تطبیق کُند می‌تواند از یادگیری
کند (\(\alpha\) پایین) یا از انتخاب پرنویز و بی‌حساس به ارزش (\(\beta\)
پایین) بیاید. جهت همبستگی به طرح بستگی دارد. علاوه بر این، مقادیر بالای
\(\beta\) ضعیف شناسایی می‌شوند، چون انتخاب‌های تقریباً قطعی برای بسیاری از
\(\beta\)های بزرگ یکسان به نظر می‌رسند.

\textbf{5.} بینایی انحراف معیار 2 و لامسه انحراف معیار 2 دارد. وزن‌های
بهینه و انحراف معیار ترکیبی پیش‌بینی‌شده چقدر است؟ \textbf{پاسخ} پایایی‌های
برابر \(w_V = w_H = 0.5\) می‌دهند، و انحراف معیار ترکیبی
\(\sqrt{4 \times 4/8} = \sqrt{2} = 1.41\) است، بهتر از هر نشانه به
تنهایی (2).

\textbf{پاسخ} علوم شناختی \ensuremath{\leftarrow} هوش مصنوعی: یادگیری خطامحور
\lr{(Rescorla\lr{--}Wagner} / قاعده‌ی دلتا) که به آموزش شبکه‌های عصبی و یادگیری
تقویتی منجر شد. هوش مصنوعی \ensuremath{\leftarrow} علوم شناختی: شبکه‌های کانولوشنی عمیق که
به‌عنوان مدل جریان بینایی شکمی \lr{(ventral visual stream)} استفاده
می‌شوند، یا عامل‌های \lr{RL} و مدل‌های زبانی که به‌عنوان مدل رفتار انسان
پیشنهاد می‌شوند.

\textbf{7.} یک پژوهشگر نشان می‌دهد پاسخ‌های یک مدل بزرگ زبانی در یک تکلیف
استدلال همبستگی بالایی با پاسخ‌های انسانی دارد. پیش از پذیرفتنش به‌عنوان
یک مدل شناختی چه شواهد دیگری می‌خواهید؟ \textbf{پاسخ} پیش‌بینی شرایط تازه
و تفاوت‌های فردی (تعمیم)؛ یک مقایسه‌ی منصفانه با مدل‌های شناختی ساده‌تر؛
آزمون‌هایی که سازوکار مدل را از جایگزین‌هایی که همان خروجی‌ها را تولید
می‌کنند تمیز دهند (تقلید)؛ و نوعی تناظر بین حالت‌های درونی‌اش و سازه‌های
روان‌شناختی.

\begin{center}\rule{0.5\linewidth}{0.5pt}\end{center}

\subsection{تکلیف}

وقت \lr{lab} این هفته را برای کار روی پروژه و جمع‌بندی درس استفاده کنید.

\begin{center}\rule{0.5\linewidth}{0.5pt}\end{center}

\subsection{مطالعه‌ی بیشتر}

\begin{itemize}
\tightlist
\item
  \textbf{\lr{Verguts} (2022)}: فصل‌های قاعده‌ی دلتا و یادگیری تقویتی
  (یادگیری خطامحور و پایه‌ی عصبی‌اش) {[}فصل‌ها بررسی شود{]}. \emph{(منبع
  اصلی برای بخش 1.)}
\item
  \textbf{\lr{Farrell} \lr{\&} \lr{Lewandowsky} (2018)}: فصل‌های مدل‌های شبکه‌ی
  عصبی و مدل‌ها در علوم اعصاب {[}شماره‌ی فصل‌ها بررسی شود{]}.
\item
  \textbf{\lr{Anderson} (2014)}: فصل‌های شبکه‌های عصبی و یادگیری، که
  مدل‌های شناختی و علوم اعصاب محاسباتی را به هم وصل می‌کنند {[}فصل‌ها بررسی
  شود{]}.
\item
  \emph{اختیاری، خارج از فهرست درس:} \lr{Ernst, M. O.,} \lr{\&} \lr{Banks, M.
  S.} (2002). \lr{Humans integrate visual and haptic information in a
  statistically optimal fashion.} \emph{\lr{Nature,} 415}, 429\lr{--}433.
  \lr{Wilson, R. C.,} \lr{\&} \lr{Collins, A. G. E.} (2019). \lr{Ten simple
  rules for the computational modeling of behavioral data.}
  \emph{\lr{eLife,} 8}, \lr{e49547} (مثال اصلی‌اش یک تکلیف \lr{bandit}
  مثل تکلیف امروز است).
\end{itemize}

\emph{پیش‌نمایش هفته‌ی 16:} ارائه‌های پروژه‌ی پایانی، و یک جلسه‌ی پایانی
درباره‌ی شیوه‌ی خوب مدل‌سازی.
